A direction vector $v$ must be perpendicular to both normal vectors $(1,-1,1)$ and $(2,3,-1)$. Their cross product is $(-2,3,5)$. It is nonzero, and its dot products with both normals vanish. Thus $\boxed{(-2,3,5)}$ is a direction vector.
\begin{asy}
import graph;
usepackage("amsmath");
import geometry;
size(7cm);
draw((-2,-1)--(2,1),Arrow);label("$v=(-2,3,5)$",(2,1),(1,1));draw((-2,1)--(2,-1)--(3,0)--(-1,2)--cycle,dashed);label("plane intersection",(0,-1.4),(0,-1));
\end{asy}

